% !TeX root = ../../Book3.tex
\section{万有覆盖}
\label{b3:sec:1103}

万有覆盖用路径类制造一张使所有闭路都能收缩的曲面。与最大解析延拓不同，这个构造只取决于底空间拓扑，随后才拉回复结构。

\subsection{万有覆盖}
\label{b3:subsec:110301}

固定 \(a\in Y\)。令 \(\widetilde Y\) 为从 \(a\) 出发的路径端点固定同伦类，投影取终点。对终点 \(y\) 的一个小坐标圆盘 \(V\)，以及一条到 \(y\) 的路径 \(\gamma\)，把 \([\gamma*\eta]\)、其中 \(\eta\) 在 \(V\) 从 \(y\) 到变动终点，作为一个坐标片。圆盘单连通使类只依赖终点，故这个片与 \(V\) 同胚。两片若有同一点，则在共同更小盘的部分相同，否则互不交。

下面构造\term[wanyou-fugai]{万有覆盖}[universal covering]，其源空间单连通。

\begin{theorem}[万有覆盖]\label{b3:thm:11-universal-cover}
每张连通 Riemann 曲面具有单连通的全纯覆盖 \(p:\widetilde Y\rightarrow Y\)，给定一个纤维点后它在覆盖同构意义下唯一。
\end{theorem}
\begin{proof}
上述片构成拓扑，投影在每片同胚，故是覆盖。路径类可沿其代表逐步提升，说明 \(\widetilde Y\) 路径连通。端点不同的两类由不交底盘分开；端点相同而类不同，则同一小盘的两层不交，所以空间 Hausdorff。底曲面有可数小盘基，有限盘链及起始类确定可达层，因此得到可数基。

给覆盖空间基点闭路 \(\beta\)，其投影为 \(\alpha\)。路径提升定义使 \(\beta\) 的终点类恰为 \([\alpha]\)；闭合意味着 \([\alpha]\) 为平凡类。把底闭路的收缩同伦提升，得到 \(\beta\) 的收缩，故源单连通。拉回坐标给全纯结构。若另有单连通覆盖，沿同一路径从选定点提升取末点；同伦提升使它依类良定义，并在每小片是底坐标的恒等对应。反向同样构造，两复合因路径提升唯一性为恒等，得到唯一同构。
\end{proof}

\subsection{基本群作用}
\label{b3:subsec:110302}

基点闭路类 \([\alpha]\in\pi_1(Y,a)\) 对路径类作 \([\gamma]\mapsto[\alpha*\gamma]\)。同伦相容性与逆路径给这是甲板变换。任意甲板变换由它对基点纤维的像决定，而该像正是一个闭路类，所以这种作用给甲板群与基本群的对应。

\ProseParagraphBreak
拼接顺序采用先走左因子再走右因子，群表示的写法随这个约定固定；如果文献用另一种路径乘法，可能出现反群。几何结论是相同的：基本群记录万有覆盖各层，闭路提升末点记录该类所对应的甲板平移。

\ProseParagraphBreak
连通覆盖的纤维置换作用是传递的。任意两点 \(x_1,x_2\) 在同一纤维，可在源曲面中选连接路径；其投影是目标闭路，从 \(x_1\) 提升后终点为 \(x_2\)。但传递性不等于存在把任意纤维点交换的甲板变换。甲板变换必须同时与整张覆盖相容，不是一个纤维上的任意置换。万有覆盖的甲板群才有正文所证的自由传递性。

\ProseParagraphBreak
例如有限覆盖的基本群子群为 \(H\)，闭路作用可以在余类上描述，甲板群则只由保持 \(H\) 的整体对称产生。如果 \(H\) 不是正规子群，一张连通覆盖可能有很少的甲板变换，尽管所有纤维点仍由闭路作用互相可达。这说明解析延拓产生的根置换，与某张曲面自身的全局对称，也须分别记录；二者在高度对称的 \(z^m\) 例子中恰好紧密对应，不应把该例的特殊性写成一般事实。

\subsection{甲板群}
\label{b3:subsec:110303}

\begin{proposition}\label{b3:prop:11-universal-deck}
万有覆盖的甲板群在每个纤维上传递且自由，底曲面与商 \(\widetilde Y/\operatorname{Deck}(p)\) 自然同构。
\end{proposition}
\begin{proof}
两个同终点路径类 \([\gamma],[\eta]\) 的差可由闭路 \(\eta*\gamma^{-1}\) 补偿，前置这条闭路将第一类变成第二类。故传递；自由性已由甲板固定点命题证明。投影每个纤维恰是一个轨道，诱导商到 \(Y\) 的双射。小层邻域上投影同胚，商映射同样局部同胚，因此双射为同胚；坐标相同使它双全纯。
\end{proof}
反向，若一组全纯自同构自由且适当地不连续作用于曲面，商具有拉回自局部不交邻域的复结构，投影为覆盖。自由性不能删除，否则商坐标将包含分歧；局部不交条件不能删除，否则商可能不再 Hausdorff。

\begin{proposition}[甲板作用的局部离散性]\label{b3:prop:11-deck-proper}
万有覆盖的甲板群在源曲面上自由作用。每一点有小邻域 \(V\)，使 \(T(V)\cap V=\varnothing\) 对所有非恒等甲板变换成立；任意两个紧集 \(K,L\) 之间只有有限个甲板变换满足 \(T(K)\cap L\ne\varnothing\)。
\end{proposition}
\begin{proof}
取均匀覆盖邻域的一层 \(V\)。若 \(T(V)\) 与 \(V\) 相交，两者作为层便相同；由于投影在该层一一对应，\(T\) 在 \(V\) 为恒等，唯一提升给全局恒等。若紧集结论失败，取不同 \(T_n\) 与 \(x_n\in K\)，使 \(T_nx_n\in L\)。子列上 \(x_n\to x\)、\(T_nx_n\to y\)，且 \(p(x)=p(y)\)。取该目标点的一个均匀邻域，令含 \(x,y\) 的层为 \(V,W\)。充分大时 \(x_n\in V\)、\(T_nx_n\in W\)，所以所有 \(T_n\) 都把 \(V\) 送到 \(W\)。投影一一性使它们在 \(V\) 相同，从而全局相同，矛盾。
\end{proof}
这给商的局部坐标：在 \(V\) 中没有不同点被甲板群识别，所以商映射在那里是同胚；目标曲面恰由这些图拼成。紧集有限相交也是后面在双曲圆盘中求最短甲板距离、证明商度量完备的重要事实。不能只说群“离散”而省略到底是在变换空间中离散还是轨道在曲面中局部离散。

\subsection{Riemann 曲面观点}
\label{b3:subsec:110304}

函数芽曲面若其投影确为覆盖，可以由底空间基本群的作用描述；一般最大芽空间的局部双全纯投影未必完整覆盖，因为某些路径遇到奇点无法继续。万有覆盖则对每条底路径都能提升。二者的关系是：若一个初始芽能沿底空间全部路径延拓，拉到万有覆盖后由单值性定理成为单值全纯函数，下降到底空间当且仅当它对全部甲板变换不变。

\ProseParagraphBreak
这给对数与幂函数一个统一位置，也为一般区域的 Poincaré 度量铺路。局部坐标随甲板变换而变，但圆盘的内在密度保持不变，所以可以下降。要知道万有覆盖是否圆盘，还需要统一化输入，不能由拓扑万有覆盖存在性直接推出。

\ProseParagraphBreak
拓扑覆盖携带复结构的过程具有唯一性。每一层用目标坐标与投影复合成源坐标，层间转换正是目标转换，故全纯。若另一复结构也使投影局部双全纯，两结构在每个层上都由同一个目标坐标拉回，必相容。这意味着覆盖问题的主要新资料在全局路径与纤维，局部解析结构已经被目标决定。分歧覆盖则还需补充局部重数，不能通过普通覆盖的坐标拉回直接越过分歧点。

\ProseParagraphBreak
读者可以把本章三个构造连在一起核验：先由路径同伦类形成万有覆盖，再以层坐标给它复结构，最后以闭路拼接给甲板变换。任一构造都不使用统一化；只有在下一章把这个已建曲面识别为球面、平面或圆盘时，才加入深层分析输入。先构造对象再识别模型，使后面度量的定义不依赖一个尚未写清的覆盖空间。
